\documentclass[10pt,a4paper]{amsart}
\usepackage[margin=19mm]{geometry}
\usepackage{amsmath,amssymb,mathtools}
\usepackage[scr=rsfs,cal=euler]{mathalfa}
\usepackage{xfrac}
\usepackage[unicode,hidelinks,bookmarks=false]{hyperref}
\newcommand{\ie}{i.e.\ }
\newcommand{\cf}{cf.\ }
\newcommand{\resp}{respectively\ }
\newcommand{\Subsection}{\subsection}
\providecommand{\nobreakdash}{\nobreak\hbox{-}}
\setlength{\emergencystretch}{2em}
\multlinegap0pt
\usepackage{algorithm}\usepackage{algpseudocode}
\usepackage{amssymb}
\usepackage{float}
\usepackage{xcolor}
\usepackage{tikz}
\usepackage{mathtools}
\usepackage{dsfont}
\newtheorem{theorem}{Theorem}
\newtheorem{proposition}{Proposition}
\title{\textbf{On the restricted Hanoi Graphs}}
\author{El-Mehdi Mehiri}
\date{Source: arXiv:2304.03857v1 (CC BY 4.0)}
\begin{document}
\maketitle

\noindent\emph{Source and licence.} \textbf{On the restricted Hanoi Graphs}. Version arXiv:2304.03857v1 (CC BY 4.0), distributed by arXiv under the Creative Commons Attribution 4.0 International licence, http://creativecommons.org/licenses/by/4.0/, as recorded in the arXiv licence metadata for this exact version and shown on the abstract page https://arxiv.org/abs/2304.03857v1. Full licence notice: \texttt{licenses/arxiv-2304.03857-CC-BY-4.0.txt}.\par\smallskip

\noindent\textit{Editorial scope.} Counted unit: the article's proposition proposition (source0.tex lines 840--842). The article's complete original proof of it is delivered with it (source0.tex lines 843--859, one \texttt{proof} environment of 17 lines). No mathematics is written, repaired or re-derived: the statement and the proof are the article's own lines, in the article's own notation, and no retained line is substituted or re-flowed.  Retained uncounted as the source-local context the proof needs: lines 772--781; lines 797--835.  Nothing was omitted from any retained span, so the package carries no omission record.  No retained line contains a citation, so the wrapper carries no bibliography and no bibliography entry.  No retained line contains a figure environment, an image-inclusion command or drawing code, so no image is delivered and the package declares no asset dependency.  Editorial formatting only, in addition: an AMS \texttt{amsart} preamble that restates the article's own declarations and macros used by the retained lines, namely the environments \texttt{theorem}, \texttt{proposition} on separate counters, together with the standard packages those lines need.  The retained environments print in this document as Theorem 1, Proposition 1 in order of appearance, whereas the article prints the same items with the numbering of its own full document.  The complete native source of the article as distributed in the arXiv e-print for this exact version is bundled unchanged as \texttt{sources/139843/source0.tex}.
\begin{theorem}
\label{theoremI} For all integers $n\geq0$, $m\geq 3$, and a digraph $D$ of $m$ vertices, we have  $(u,v)\in A(H_{n}^{m}(D))$ only if there exists a unique  $ k \in \mathcal{N}$ such that the following conditions are fulfilled :
\begin{enumerate}
    \item $u_{k}\neq v_{k}$ and  $u_{d}=v_{d}$ for all $d\in N\setminus\{k\}$; (only one disc can be moved)
    \item $u_{d}\neq u_{k}$  for all $d<k$; (the disc to be moved is a topmost disc) 
    \item  $v_{d}\neq v_{k}$ for all $d<k$; (a disc cannot reside on a smaller one) 
    \item $(u_{k},v_{k})\in A(D)$; (discs are allowed to move from peg $u_{k}$ to peg $v_{k}$)
\end{enumerate}
    
\end{theorem}

Theorem \ref{theoremI} allows us to design the simple Algorithm \ref{alg:gen} to generate the restricted Hanoi graphs  $H_{n}^{m}(D)$ for any $n\geq 0$, $m\geq 3$, and digraph of movements $D$ using the boolean function described in Algorithm \ref{alg:cond} that returns true if and only if a candidate arc $(u,v)\in V(H_{n}^{m}(D))$ is admissible to be in  $A(H_{n}^{m}(D))$. 
\begin{algorithm}[H]
		\caption{A boolean function to decide whether an arc is in $H_{n}^{m}(D)$ or not }\label{alg:cond}
		\begin{algorithmic}[1]
			\Function{Condition}{$u,v$}
                \State $\mathcal{K}=\{d\in\mathcal{N}\mid u_{d}\neq v_{d}\};$
                \If{$|\mathcal{K}|\neq 1$}
                \State \Return false;
                \EndIf
                \State let $k$ be the  element of the singleton set $\mathcal{K}$;
                \If{$(u_{k},v_{k})\notin A(D)$}
                \State \Return false;
                \EndIf
                \If{$k\neq 1$}
                \For{$d$ from $k-1$ to $1$}
                \If{$u_{d}=u_{k}$ or $v_{d}=v_{k}$}
                \State \Return false;
                \EndIf
                \EndFor
                \EndIf
                \State \Return true;
			\EndFunction
		\end{algorithmic}
	\end{algorithm}
	\begin{algorithm}[H]
		\caption{The restricted Hanoi graphs generator }\label{alg:gen}
		\begin{algorithmic}[1]
			\Require $n$, $m$ and $D=(\mathcal{M},A(D))$;
                \Ensure The restricted Hanoi graph $H_{n}^{m}(D)$;
                \State generate all vertices of $V(H_{n}^{m}(D))$;
                \State  $A(H_{n}^{m}(D)):=\emptyset$;
                \For{$u\neq v\in V(H_{n}^{m}(D))$}
                \If{$\textsc{Condition}(u,v)=true$} 
                \State $A(H_{n}^{m}(D)):=A(H_{n}^{m}(D))\cup\{(u,v)\}$;
                \EndIf
                \EndFor 
                \State \Return $H_{n}^{m}(D)$;
		\end{algorithmic}
	\end{algorithm}

\begin{proposition}
     For all integers $n\geq0$, $m\geq 3$, and a digraph $D$ of $m$ vertices, the restricted Hanoi graph $H_{n}^{m}(D)$ can be generated  in a time complexity $\mathcal{O}(m^{2n}(m^{2}+n))$.
\end{proposition}

\begin{proof}
    Consider Algorithm \ref{alg:gen} as a generator of  restricted Hanoi graph $H_{n}^{m}(D)$, and let $f(n,m)$ and $g(n,m)$ be the number of operations needed to perform Algorithm \ref{alg:gen} and \ref{alg:cond} respectively, then
    \begin{align*}
        f(n,m)&=m^{n}+m^{n}(m^{n}-1)g(n,m)+1.
    \end{align*}
   In the worst case where $D=\overleftrightarrow{K}_{m}$ and $k\neq 1$, we would have
    \begin{align*}
        g(n,m)&=m^{2}-m+3n-1\\
        &=\mathcal{O}(m^{2}+n).
    \end{align*}
    Thus, 
    \begin{align*}
        \mathcal{O}(f(n,m))&=m^{n}+(m^{2n}-m^{n})\mathcal{O}(m^{2}+n)\\
        &=\mathcal{O}(m^{2n}(m^{2}+n)).
    \end{align*}
    
\end{proof}

\end{document}
